\section{Tools resources prompts and skills remain distinct}\label{sec:C19:primitives}
MCP exposes several kinds of objects whose roles should remain separate. Tools are invocable capabilities with schemas; resources expose addressable context; prompts provide reusable interaction templates~\cite{mcp2026spec}. A host decides how those objects enter the model's working view. The protocol does not specify a universal context-selection, memory-consolidation, or reasoning policy. A shared transport therefore does not make every object an instruction or every server a memory system.

The current tool interface separates listing from invocation through \path{tools/list} and \path{tools/call}. Available tools can depend on request authorization, and descriptions and schemas guide the caller~\cite{mcp2026tools}. The host still needs grounding, argument validation, and an authority check before effects. A tool annotation is not a substitute for an enforced permission boundary. Chapter~\ref{ch:C18} analyzes these application-level contracts.

Resources have URI identities and can be listed, read, or described through templates~\cite{mcp2026resources}. A resource URI is a reference in the server's namespace, not necessarily a local filesystem path accessible to the client. The host should preserve the server and version context when storing or citing retrieved content. If a resource later changes, a previous citation requires a content snapshot or version reference to remain reproducible. A stable URI alone does not imply immutable evidence.

A skill is procedural knowledge packaged for selective use, as developed in Chapter~\ref{ch:C15}. It may be delivered through a protocol extension or represented using resources, prompts, and tools, but those representations do not collapse its meaning into the underlying primitive. The Agent Skills specification concerns the skill artifact; MCP concerns exchange and invocation; the model policy decides applicability. A skill can instruct the use of a tool without having authority to grant access to it~\cite{agentskills2026spec}.

This separation has a concrete security consequence. A retrieved resource can contain adversarial text that asks the agent to change its goal. The host should label it as untrusted evidence rather than promote it to a privileged instruction. A prompt template can be selected by an authorized user while still containing external fields that require careful interpolation. A tool result can contain both structured data and explanatory prose with different trust properties. These distinctions must survive serialization and compaction.

The proposed implementation stores a typed provenance envelope for every imported object: server identity, object kind, object identifier, schema or content version, access scope, and insertion role. A context assembler then decides what is visible without losing those labels. This is an application design, not a protocol guarantee. It makes cross-layer errors observable when the same bytes are later retrieved as memory or passed to another agent.

